\section{领域一：Agentic 架构与编排（27\%）}

这是权重最高的领域，考察候选人设计和实现 agent 系统的能力，包括 agentic loop 机制、子 Agent 编排以及安全关键场景下的 hook 设计。

\subsection{三大反模式}

考试会直接测试你能否识别以下三个常见的反模式：

\begin{warningbox}{必须立即拒绝的三种 Agent 停止策略}
\begin{enumerate}[leftmargin=1.5em]
    \item \textbf{解析自然语言判断循环终止}：通过分析模型回复的文本内容来决定是否停止循环。这种方式不可靠，因为自然语言具有歧义性。
    \item \textbf{任意迭代次数上限作为主要停止机制}：例如"最多运行 10 轮"。这不是一个语义上合理的终止条件，应基于 \texttt{stop\_reason} 来判断。
    \item \textbf{检查 assistant 文本作为完成指示器}：模型输出文本并不意味着任务完成，正确做法是检查 API 返回的 \texttt{stop\_reason} 字段。
\end{enumerate}
\end{warningbox}

\subsection{子 Agent 的上下文隔离}

\begin{importantbox}{子 Agent 不共享 Coordinator 的内存}
这是原文作者强调的"最大错误"：很多人假设子 Agent（subagent）可以访问 coordinator 的上下文。实际上，\textbf{子 Agent 的上下文是完全隔离的}，所有信息必须通过 prompt 显式传递。
\end{importantbox}

这意味着在设计多 Agent 系统时，你需要：
\begin{itemize}[leftmargin=1.5em]
    \item 明确定义 coordinator 向每个 subagent 传递的信息边界
    \item 在 subagent 的 prompt 中包含所有必要的上下文
    \item 设计清晰的返回值结构，确保 coordinator 能整合各 subagent 的结果
\end{itemize}

\subsection{安全关键场景的 Hook 机制}

\begin{importantbox}{高风险场景必须用程序化 Hook 而非 Prompt 指令}
当涉及金融交易或安全敏感操作时，仅靠 prompt 中的指令是不够的。必须通过 \textbf{hook} 和 \textbf{prerequisite gate} 以编程方式强制执行工具调用顺序。
\end{importantbox}

例如，在一个涉及退款操作的 Agent 中：
\begin{itemize}[leftmargin=1.5em]
    \item 使用 \texttt{PostToolUse} hook 规范化数据格式
    \item 使用工具调用拦截 hook 阻止违反策略的操作
    \item 通过 prerequisite gate 确保在执行退款前必须先完成身份验证
\end{itemize}

\subsection{推荐实践项目}

\begin{knowledgebox}{动手练习：多工具 Agent}
构建一个包含 3--4 个 MCP 工具的 Agent，实现以下功能：
\begin{itemize}[leftmargin=1.5em]
    \item 正确的 \texttt{stop\_reason} 处理
    \item 一个 \texttt{PostToolUse} hook 用于规范化数据格式
    \item 一个工具调用拦截 hook 用于阻止策略违规
\end{itemize}
这一个练习就能覆盖领域一的大部分考点。
\end{knowledgebox}

\subsection{本章小结}

Agentic 架构的核心在于三点：正确使用 \texttt{stop\_reason} 而非自然语言解析来控制 Agent 循环；理解子 Agent 上下文隔离的本质并显式传递信息；在高风险场景中用程序化 hook 而非 prompt 指令来保障安全。

